\section{Introduction}
\label{sec:introduction}

Automated medical image segmentation is a cornerstone of modern
computer-aided diagnosis, enabling precise delineation of anatomical
structures and pathological regions in MRI, CT, and ultrasound scans.
Deep learning approaches---particularly fully-convolutional
UNet-style architectures~\cite{ronneberger2015unet}---have achieved
near-radiologist accuracy on a growing number of benchmarks.
Yet their clinical deployment remains constrained by a fundamental
tension: state-of-the-art models are data-hungry, whereas patient data
are sensitive and subject to strict privacy legislation (e.g.,
GDPR, HIPAA).

\paragraph{The federated learning promise.}
McMahan \emph{et al.}~\cite{mcmahan2017fedavg} proposed \emph{Federated
Learning} (FL) as a paradigm shift: rather than aggregating raw images on a
central server, model gradients (or weights) are communicated, keeping
patient data on-premise.
FL has since attracted significant interest for medical imaging
applications~\cite{sheller2018federated,nvidia2019clara,bercea2022federated,liu2024fedfms,skorupko2025fednnunet},
but two important challenges remain largely unresolved:

\begin{enumerate}
  \item \textbf{Statistical heterogeneity.}
    Clinical data across institutions follow non-identical distributions
    (non-IID), caused by scanner variability, patient demographics, and
    annotation protocols.
    Standard FedAvg degrades substantially under high heterogeneity.
  \item \textbf{Annotation scarcity.}
    Pixel-level segmentation masks require expert radiologists and are
    expensive to obtain.
    In FL settings each client can only annotate its own local pool,
    amplifying the labelling bottleneck.
\end{enumerate}

\paragraph{Contributions.}
We address both challenges with the following contributions:
\begin{itemize}
  \item A unified FL framework for 2-D medical image segmentation
    supporting four aggregation algorithms: FedAvg~\cite{mcmahan2017fedavg},
    FedProx~\cite{li2020fedprox}, FedNova~\cite{wang2020fednova}, and
    FedPer~\cite{arivazhagan2019fedper}.
  \item Systematic evaluation under Dirichlet non-IID partitioning on four
    clinically diverse datasets: ACDC, BUSI, FUGC, and TN3K.
  \item \textbf{FedAL}: a federated active-learning protocol that integrates
    entropy-based query selection into the FL communication rounds, reducing
    labelling cost while maintaining competitive segmentation performance.
  \item A 1\textsuperscript{st}-place result on the FUGC 2025 challenge,
    demonstrating the practical effectiveness of the proposed framework.
\end{itemize}
